\section{Related Work}
\label{sec:related}

\textbf{Counterfactual reliance.} The Visual Necessity Score of \emph{To See or To Please}~\citep{hong2026see} is the closest prior instrument: it applies full, blind, noise, and conflict conditions to nine VLMs over 9,000 model--sample pairs and reports family scaling and a training-free selective-prediction gain. Our reliance conditions are in the same family, but our question is different: not ``what is the score'' but ``when is the score identifiable.'' That paper's no-image contrast is reported without a real-vs-real null; we show that this omission can reverse the ranking (Section~\ref{sec:results-null}). \emph{Do VLMs See or Guess?}~\citep{singla2026guess} holds the image fixed and varies the question, isolating textual-prior reliance. Our result is complementary: when the image is removed or degraded, a collapsed or templated model can produce a large apparent contrast with no evidence extraction, so phrasing controls alone do not identify image use.

\textbf{Corruption robustness.} Medical VLM robustness is usually summarized by clean error and mean corruption error~\citep{imam2025corruption}, or by accuracy and artefact-detection rate~\citep{cheng2025robustness}. These metrics aggregate over labels and therefore cannot see an operating-point shift: when true positives and false positives rise together, accuracy can be flat or improve. We provide an explicit empirical analysis of this failure mode in chest X-ray VLM evaluation, showing that a degradation can \emph{raise} recall while precision and false positives move against robustness, and that both must be reported to detect it. Stress-testing suites report comparable fragility under prompt and workflow changes~\citep{gu2025illusion}, and grounding studies question whether medical VLMs use image evidence at all~\citep{felizzi2025grounded}.

\textbf{Conservative default-normal behaviour.} \emph{Reliability Stress Tests}~\citep{yang2026stress} shows that exact-match accuracy overstates conservative models that default to ``Normal,'' using CheXagent~\citep{chen2024chexagent} and MedGemma~\citep{sellergren2025medgemma} under prompt and workflow variation. We replicate the phenomenon at larger scale and add the contract axis: the \emph{form} of the shortcut (confident normal vs.\ abstention) is decided by the answer schema, not by the model (Section~\ref{sec:results-schema}). Our label-level analysis also separates ``default normal'' from ``collapsed output,'' which that framing does not distinguish.

\textbf{Abstention and refusal.} \emph{Knowing When Not to Answer}~\citep{madhusudhan2026knowing} evaluates abstention on unanswerable multimodal inputs and shows that benchmarks assuming answerability miss the failure mode. Our blank-control analysis shows a complementary artifact: with a prose contract, abstention on a blank image is common and produces a spurious ``strong reliance'' reading unless the blank template is inspected (Section~\ref{sec:results-null}).

\textbf{Reasoning length and termination.} A separate line of work shows that longer reasoning is not monotonically better: accuracy follows an inverted-U with chain-of-thought length~\citep{wu2026whenmore}, extended thinking can abandon answers the model already had~\citep{zhou2026overthinking}, and truncated traces can hurt more than no reasoning~\citep{su2026broken}. That literature measures accuracy under a fixed budget. Our second-model result is a measurement consequence: when a reasoning checkpoint is asked for structured output, non-termination removes the cell from the denominator, so coverage itself becomes a function of the corruption (Section~\ref{sec:results-second}).

\textbf{Calibration and selective prediction.} Calibration is conventionally measured with ECE and corrected by temperature scaling~\citep{guo2017calibration}, and its fragility under distribution shift is well documented~\citep{ovadia2019trust}. Risk--coverage analysis is standard in medical machine learning~\citep{geifman2017selective,oakdenrayner2021limits}. Our contribution is negative and specific: under a pre-declared cost model, the composite policy never accepts at the image level (accept rate 0.000), while at the finding level it accepts 17--41\% of items but with an accepted error rate of 5.7--21.7\%, above the 5\% budget. Coverage at 5\%/10\%/20\% risk is 0.000 at the image level; at the finding level it reaches at most 0.093. Adding shift or reliance signals to confidence moves finding-level AUROC by at most 0.034, so better error ranking---where it exists---does not translate into a usable decision rule (Section~\ref{sec:results-selective}).

\textbf{Report generation and hallucination.} Chest X-ray report generation is evaluated by factual completeness and hallucination as well as fluency~\citep{heiman2025factchexcker}. Our Open-i layer applies a deliberately simple labeller to both the reference and the generated report; a CheXbert cross-check shows that the omission ordering is labeller-dependent while the hallucination ordering is preserved in direction, so we report only the latter as stable (Section~\ref{sec:results-report}).

\textbf{Datasets and label provenance.} The public CXR corpora used here were built for weak supervision and report-derived labels~\citep{wang2017chestxray8,bustos2020padchest,irvin2019chexpert,johnson2019mimic,chambon2024chexpertplus}. Their label noise, site composition, and evaluation conventions are known to bound what a benchmark can claim~\citep{rafferty2026limitations}, and confounding between pathology and site can survive a random split~\citep{zech2018confounding}. Supervised CXR classifiers (DenseNet, CheXNet) remain the reference point for discrimination~\citep{rajpurkar2017chexnet,huang2017densely}. To separate model weakness from intrinsic task difficulty, we train a patient-disjoint DenseNet-121 baseline alongside the VLMs (macro AUROC 0.733; Section~\ref{sec:results-baseline}). We therefore treat label provenance and case mix as measurement declarations rather than as caveats added after the fact (Table~\ref{tab:prevalence}).

\textbf{Medical foundation models and reporting standards.} Generalist medical AI and medical LLMs are motivated by scale and breadth~\citep{moor2023foundation,singhal2023medpalm,bommasani2021foundation}, but benchmark claims require documented data and model provenance. Datasheets~\citep{gebru2021datasheets} and model cards~\citep{mitchell2019model} are the established instruments for that documentation; our four preconditions are an analogous, task-specific contract for an evaluation number.

\textbf{Measurement validity and direct competitors.} Several 2026 works ask adjacent questions (Table~\ref{tab:competitors}). \emph{Answer, Clarify, or Abstain}~\citep{nguyen2026answer} proposes fine-grained selective prediction for medical VLMs, and CALM-VLM~\citep{dhinagar2026calm} combines calibration with selective prediction for brain MRI; both act at the output level but do not audit the measurement preconditions. \emph{Beyond Accuracy}~\citep{zafar2026beyond} evaluates visual grounding in multimodal medical reasoning, and Buckley et al.~\citep{buckley2026exploit} show that multimodal foundation models can exploit text shortcuts in medical prediction. Qu et al.~\citep{qu2026accuracy} find in mammography that answer accuracy overstates evidence grounding and abstention reliability; our CXR multi-label protocol adds the responsiveness, real-vs-real null, operating-point, and contract-coverage audits, and reports the selective-trust negative result. Two negative results are also close: Fisher~\citep{fisher2025negative} reports that vision-language foundation models may fail to transfer to CXR classification, and AT-CXR~\citep{li2025atcxr} proposes uncertainty-aware triage.

\begin{table}[t]
\centering
\footnotesize
\setlength{\tabcolsep}{3pt}
\caption{Positioning against the closest 2025--2026 work. ``Overlap'' is what the prior work already establishes; ``increment'' is what this paper adds beyond it. No single component is claimed as new.}
\label{tab:competitors}
\begin{tabular}{@{}>{\raggedright\arraybackslash}p{0.19\textwidth}>{\raggedright\arraybackslash}p{0.22\textwidth}>{\raggedright\arraybackslash}p{0.25\textwidth}>{\raggedright\arraybackslash}p{0.24\textwidth}@{}}
\toprule
Work & Main question & Overlap & This paper's increment \\
\midrule
\emph{To See or To Please}~\citep{hong2026see} & VLM reliance under blind/noise/conflict & reliance score (VRS), family scaling & real-vs-real null; responsiveness; operating point; contract coverage \\
\emph{Reliability Stress Tests}~\citep{yang2026stress} & prompt/workflow stress on CXR VLMs & exact match overstates default-Normal & contract axis; three layers; label-level and selective metrics \\
\emph{Do VLMs See or Guess?}~\citep{singla2026guess} & phrasing-controlled textual prior & no-image contrast & collapsed-output/template null; responsiveness audit \\
\emph{Answer, Clarify, or Abstain}~\citep{nguyen2026answer} & fine-grained selective policy for medical VLMs & answer/clarify/abstain decisions & negative result; four preconditions before a policy is trusted \\
\emph{Accuracy Overstates}~\citep{qu2026accuracy} & mammography grounding and abstention & accuracy overstates reliability & CXR multi-label; cross-dataset layer; operating point \\
CALM-VLM~\citep{dhinagar2026calm} & calibration + selective prediction, brain MRI & calibration and selective metrics & CXR; contract coverage; cross-dataset collapse \\
AT-CXR~\citep{li2025atcxr} & uncertainty-aware agentic CXR triage & uncertainty for triage & audits the measurement, not the triage policy \\
VL FMs do not transfer~\citep{fisher2025negative} & negative transfer to CXR classification & transfer failure & mechanism: responsiveness and null model; three layers \\
\bottomrule
\end{tabular}
\end{table}

Our increment is not a new score or policy. It is the joint audit of responsiveness, null model, operating point, and contract coverage, plus the finding-level negative result that the composite policy accepts 17--41\% of items but at 5.7--21.7\% accepted error, above a 5\% budget. We do not claim priority for any single component.

\textbf{Concurrent work.} Seven 2026 papers overlap directly. Nguyen et al.~\citep{nguyen2026answer} propose a fine-grained answer/clarify/abstain policy with a condition-aware reliability layer; our negative result is the precondition audit that must precede such a policy, and our 5\% finding-level budget is deliberately stricter than their triage setting. Lotfinia et al.~\citep{lotfinia2026needimage} audit nine chest-radiograph systems with region occlusion, irrelevant-region occlusion and same-label swaps, and classify systems as image-ignoring, unstable or selectively image-using. Sadanandan and Behzadan~\citep{sadanandan2026consistent} cross consistency under paraphrase with image removal into a four-quadrant safety taxonomy. Akhtar et al.~\citep{akhtar2026beyond} audit cohort, prompt, negative-spectrum, prevalence and threshold changes over 12,200 chest radiographs and find that prompt and cohort shifts move AUROC. Hao and Floresca~\citep{hao2026evidence} track baseline-correct pairs under progressive degradation and find persistent commitment beyond an evidence boundary. Kale and Bajwa~\citep{kale2026aggregate} show that aggregate accuracy hides complete diagnosis-specific failure. Khanmohammadi et al.~\citep{khanmohammadi2026calibrated} evaluate nine confidence estimators under Learn-then-Test and show that local labels buy certified coverage while estimators trained on natural images transfer poorly. Our overlap is the observation that image reliance is often unidentifiable. Our increment is the four audit preconditions, the finding-level selective negative result under three confidence definitions, and the four-model cross-modal replication. We do not claim priority for the observation; we do not evaluate local-label certification, and we do not train a correctness probe. Those are scope limits, not passed tests.

\textbf{Datasets and models.} We use NIH ChestX-ray14~\citep{wang2017chestxray8}, PadChest~\citep{bustos2020padchest}, and Open-i~\citep{demnerfushman2016openi}, with Qwen2.5-VL~\citep{bai2025qwen25vl}, Qwen3-VL~\citep{qwen2025qwen3vl}, InternVL3.5~\citep{opengvlab2025internvl35}, MedGemma~\citep{sellergren2025medgemma}, and CheXagent~\citep{chen2024chexagent}. LLaVA-Med~\citep{li2024llavamed} is run separately and excluded from quantitative comparisons because it does not obey the structured contract.
